\section{Preprocesamiento}\label{sec:preprocesamiento}

El preprocesamiento se ajusta únicamente en la partición de fit correspondiente. Normaliza categorías, crea una categoría \texttt{DESCONOCIDO}, asigna frecuencia cero a códigos de vía no vistos, estandariza continuas y genera codificación \emph{one-hot}. Mes, día semanal y hora se representan cíclicamente. Se incluyen interacciones predeclaradas de contexto vial.

El contrato persistido contiene 21 campos crudos y 169 columnas \texttt{float32} en orden fijo. Se excluyen fallecidos, lesionados, vehículos dañados, causas investigadas, identificadores y toda característica PERSONAS. El esquema, escalador, encoders, modelo, calibrador y umbrales se verifican por hash antes de inferencia.

Ajustar el preprocesador únicamente en la partición de fit es la salvaguarda que evita la fuga más común en series temporales: si la media de estandarización, la mediana de imputación o el mapa de frecuencias de vía se calcularan sobre todo el conjunto, cada partición posterior quedaría contaminada con información de su propio futuro. El desempeño resultante de este contrato se documenta en la sección~\ref{sec:modelo}.
